
\section{\BookRecordIndex{T0088}殆复结构}\label{sec:C-7}

如第 2.1 节所述，一般的殆复流形 $(M,J)$ 并不存在既无挠又与 $J$ \BookRecordIndex{T0348}相容的联络；这里\BookRecordIndex{T0348}相容是指\BookRecordIndex{T0263}平行移动与 $J$ 交换。不过，它有自然的、与 $J$ \BookRecordIndex{T0348}相容的联络 $\widehat\nabla$，其\BookRecordIndex{T0248}挠率 $T$ 是 \BookRecordIndex{T0042}Nijenhuis 张量的一个倍数。本节将推导 $\widehat\nabla$ 的公式，并证明当 $J$ 与 $\omega$ \BookRecordIndex{T0348}相容时，它化为第 3.1 节考虑的联络
\[
  \widetilde\nabla_XY:=\nabla_XY-\frac12J(\nabla_XJ)Y.
\]
我们还将建立注~\ref{rem:3-1-3} 中关于 $D_u$ 的、以 $\widehat\nabla$ 表示的显式公式。

设 $(M,J)$ 为殆复流形。假定 $M$ 配有一个被 $J$ 保持的黎曼度量，以 $\nabla$ 表示该度量的 \BookRecordIndex{T0034}Levi--Civita 联络。下面的引理表明，$\nabla J=0$ 当且仅当 $M$ 是\BookRecordIndex{T0207}凯勒流形，即 $J$ 可积且 $2$-形式 $\omega:=\langle J\cdot,\cdot\rangle$ 为闭形式。

\begin{lemma}\label{lem:C-7-1}
设 $(M,J)$ 为配有一个被 $J$ 保持的黎曼度量的殆复流形，并令 $\omega:=\langle J\cdot,\cdot\rangle\in\Omega^2(M)$。则
\begin{equation}\label{eq:C-7-1}
  (\nabla_XJ)J+J(\nabla_XJ)=0,\qquad
  \langle(\nabla_XJ)Y,Z\rangle+\langle Y,(\nabla_XJ)Z\rangle=0,
\end{equation}
\begin{equation}\label{eq:C-7-2}
  d\omega(X,Y,Z)=\langle(\nabla_XJ)Y,Z\rangle
  +\langle(\nabla_YJ)Z,X\rangle
  +\langle(\nabla_ZJ)X,Y\rangle,
\end{equation}
\begin{equation}\label{eq:C-7-3}
  N(X,Y)=\bigl(J\nabla_YJ-\nabla_{JY}J\bigr)X
    -\bigl(J\nabla_XJ-\nabla_{JX}J\bigr)Y
\end{equation}
对 $X,Y,Z\in\Vect(M)$ 成立。若 $\omega$ 为闭形式，则
\begin{equation}\label{eq:C-7-4}
  \langle X,N(Y,Z)\rangle=2\langle J(\nabla_XJ)Y,Z\rangle,
\end{equation}
\begin{equation}\label{eq:C-7-5}
  J(\nabla_{JX}J)=\nabla_XJ,\qquad
  N(X,Y)=2J(\nabla_YJ)X-2J(\nabla_XJ)Y,
\end{equation}
\begin{equation}\label{eq:C-7-6}
  \langle X,N(Y,Z)\rangle+\langle Y,N(Z,X)\rangle
    +\langle Z,N(X,Y)\rangle=0
\end{equation}
对 $X,Y,Z\in\Vect(M)$ 成立。
\end{lemma}
\begin{proof}
对恒等式 $\langle Y,JZ\rangle+\langle JY,Z\rangle=0$ 与 $J^2=-\mathrm{id}$ 求导，可得~\eqref{eq:C-7-1}。为证明~\eqref{eq:C-7-2}，取向量场 $X,Y,Z$，使形如 $\nabla_XY$ 的六个协变导数在给定点 $x\in M$ 处全为零。于是三个李括号 $[X,Y]$、$[Y,Z]$、$[Z,X]$ 在 $x$ 处均为零，并且 $\nabla_X(\omega(Y,Z))=\langle(\nabla_XJ)Y,Z\rangle$ 在 $x$ 处成立。因此在该点有
\begin{align*}
  d\omega(X,Y,Z)
  &=\nabla_X\bigl(\omega(Y,Z)\bigr)+\nabla_Y\bigl(\omega(Z,X)\bigr)
    +\nabla_Z\bigl(\omega(X,Y)\bigr)\\
  &\quad+\omega([X,Y],Z)+\omega([Y,Z],X)+\omega([Z,X],Y)\\
  &=\langle(\nabla_XJ)Y,Z\rangle+\langle(\nabla_YJ)Z,X\rangle
    +\langle(\nabla_ZJ)X,Y\rangle.
\end{align*}
这证明了~\eqref{eq:C-7-2}。关于 \BookRecordIndex{T0042}Nijenhuis 张量的公式~\eqref{eq:C-7-3} 直接由
\[
  N(X,Y):=[X,Y]+J[JX,Y]+J[X,JY]-[JX,JY]
\]
以及公式 $[X,Y]=\nabla_YX-\nabla_XY$ \BookRecordIndex{T0326}推出。（我们的符号约定见第 2.1 节。）结合~\eqref{eq:C-7-3} 与 $d\omega=0$ 时的~\eqref{eq:C-7-2}，得到
\begin{align*}
  \langle X,N(Y,Z)\rangle
  &=\langle X,J(\nabla_ZJ)Y\rangle-\langle X,J(\nabla_YJ)Z\rangle\\
  &\quad-\langle X,(\nabla_{JZ}J)Y\rangle
    +\langle X,(\nabla_{JY}J)Z\rangle\\
  &=\langle X,J(\nabla_ZJ)Y\rangle-\langle X,J(\nabla_YJ)Z\rangle\\
  &\quad+\langle JZ,(\nabla_YJ)X\rangle+\langle Y,(\nabla_XJ)JZ\rangle\\
  &\quad-\langle JY,(\nabla_ZJ)X\rangle-\langle Z,(\nabla_XJ)JY\rangle\\
  &=2\langle J(\nabla_XJ)Y,Z\rangle.
\end{align*}
最后一个等号由~\eqref{eq:C-7-1} 得到。这便证明了~\eqref{eq:C-7-4}。再有
\begin{align*}
  2\langle J(\nabla_{JX}J)Y,Z\rangle
  &=\langle JX,N(Y,Z)\rangle\\
  &=-\langle X,JN(Y,Z)\rangle\\
  &=\langle X,N(Y,JZ)\rangle\\
  &=2\langle J(\nabla_XJ)Y,JZ\rangle\\
  &=2\langle(\nabla_XJ)Y,Z\rangle.
\end{align*}
其中第一、第四个等号由~\eqref{eq:C-7-4} 得到，第三个等号来自 \BookRecordIndex{T0042}Nijenhuis 张量在两个变量上均复反线性。这证明了 $J\nabla_{JX}J=\nabla_XJ$，于是~\eqref{eq:C-7-5} 由~\eqref{eq:C-7-3} \BookRecordIndex{T0326}推出。为证明~\eqref{eq:C-7-6}，注意
\begin{align*}
  &\langle X,N(Y,Z)\rangle+\langle Y,N(Z,X)\rangle
    +\langle Z,N(X,Y)\rangle\\
  &\quad=2\langle J(\nabla_XJ)Y,Z\rangle
    +2\langle J(\nabla_YJ)Z,X\rangle
    +2\langle J(\nabla_ZJ)X,Y\rangle\\
  &\quad=-2\langle(\nabla_{JX}J)Y,Z\rangle
    -2\langle(\nabla_YJ)Z,JX\rangle
    -2\langle(\nabla_ZJ)JX,Y\rangle\\
  &\quad=0.
\end{align*}
第一个等号来自~\eqref{eq:C-7-4}，第二个来自~\eqref{eq:C-7-5}，最后一个来自 $d\omega=0$ 时的~\eqref{eq:C-7-2}。
\end{proof}

现在假设\BookRecordIndex{T0363}辛形式 $\omega$ \BookRecordIndex{T0372}驯服 $J$，并如~\eqref{eq:2-1-1} 那样取黎曼度量
\begin{equation}\label{eq:C-7-7}
  \langle v,w\rangle:=\frac12\bigl(\omega(v,Jw)+\omega(w,Jv)\bigr).
\end{equation}
我们强调，除非 $J$ 与 $\omega$ \BookRecordIndex{T0348}相容，$2$-形式 $\langle J\cdot,\cdot\rangle$ 并不等于 $\omega$。令 $\nabla$ 为度量~\eqref{eq:C-7-7} 的 \BookRecordIndex{T0034}Levi--Civita 联络。则公式
\[
  \widetilde\nabla_XY:=\nabla_XY-\frac12J(\nabla_XJ)Y
\]
定义 $TM$ 上保持 $J$ 的联络。公式~\eqref{eq:C-7-1} 表明 $\widetilde\nabla$ 还保持度量。然而，这个联络一般并非无挠的。若 $J$ 与 $\omega$ \BookRecordIndex{T0348}相容，则公式~\eqref{eq:C-7-5} 表明 $\widetilde\nabla$ 的\BookRecordIndex{T0248}挠率是 \BookRecordIndex{T0042}Nijenhuis 张量的负四分之一：
\begin{align*}
  T^{\widetilde\nabla}(X,Y)
  &:=\widetilde\nabla_YX-\widetilde\nabla_XY-[X,Y]\\
  &=\frac12J(\nabla_XJ)Y-\frac12J(\nabla_YJ)X\\
  &=-\frac14N(X,Y).
\end{align*}

若 $J$ 仅被 $\omega$ \BookRecordIndex{T0372}驯服，则宜在 $TM$ 上引入联络 $\widehat\nabla$：
\begin{equation}\label{eq:C-7-8}
  \widehat\nabla_XY
  :=\widetilde\nabla_XY-\frac14\bigl(\nabla_{JY}J+J\nabla_YJ\bigr)X
  =\nabla_XY-Q(X,Y),
\end{equation}
其中
\begin{equation}\label{eq:C-7-9}
  4Q(X,Y):=(\nabla_{JY}J)X+J(\nabla_YJ)X+2J(\nabla_XJ)Y.
\end{equation}
联络 $\widehat\nabla$ 一般不保持度量；但下一个引理表明它保持 $J$，且其\BookRecordIndex{T0248}挠率恒等于 \BookRecordIndex{T0042}Nijenhuis 张量的负四分之一。

\begin{lemma}\label{lem:C-7-2}
设 $(M,\omega)$ 为\BookRecordIndex{T0358}辛流形，$J\in\mathcal J_\tau(M,\omega)$。以 $\nabla$ 表示度量~\eqref{eq:C-7-7} 的 \BookRecordIndex{T0034}Levi--Civita 联络，并由~\eqref{eq:C-7-8} 定义 $\widehat\nabla$。则 $\widehat\nabla$ 保持\BookRecordIndex{T0088}殆复结构 $J$，且其\BookRecordIndex{T0248}挠率为
\[
  T^{\widehat\nabla}(X,Y)
  :=\widehat\nabla_YX-\widehat\nabla_XY-[X,Y]
  =-\frac14N(X,Y).
\]
若 $J$ 与 $\omega$ \BookRecordIndex{T0348}相容，则 $\widehat\nabla=\widetilde\nabla$，因此 $\widehat\nabla$ 也保持度量。
\end{lemma}
\begin{proof}
$\widehat\nabla$ 保持 $J$ 可由直接计算得出。为计算其\BookRecordIndex{T0248}挠率，注意
\begin{align*}
  4\bigl(Q(X,Y)-Q(Y,X)\bigr)
  &=(\nabla_{JY}J)X-J(\nabla_YJ)X
    -(\nabla_{JX}J)Y+J(\nabla_XJ)Y\\
  &=-N(X,Y).
\end{align*}
最后一个等号来自~\eqref{eq:C-7-3}。\sourcefootnote{原书在此将最后一个等号归于 (C.7.5)，但本引理只假设 $J$ 被 $\omega$ \BookRecordIndex{T0372}驯服；(C.7.5) 要求 $J$ 与 $\omega$ \BookRecordIndex{T0348}相容。此等号直接由无须\BookRecordIndex{T0348}相容假设的 (C.7.3) 得出。}由于 $\nabla$ 无挠，$\widehat\nabla$ 的\BookRecordIndex{T0248}挠率公式随即得到。在 $\omega$-\BookRecordIndex{T0348}相容的情形，~\eqref{eq:C-7-5} 给出 $\widehat\nabla=\widetilde\nabla$。引理得证。
\end{proof}

现在证明注~\ref{rem:3-1-3} 所述的 $D_u$ 公式。

\begin{lemma}\label{lem:C-7-3}
若 $J$ 被 $\omega$ \BookRecordIndex{T0372}驯服，且 $u:\Sigma\to M$ 是光滑映射，则
\begin{equation}\label{eq:C-7-10}
  D_u\xi=(\widehat\nabla\xi)^{0,1}
    +\frac14N_J\bigl(\xi,\partial_J(u)\bigr)
    +\frac14\bigl(J(\nabla_{\bar\partial_J(u)}J)
                     +\nabla_{J\bar\partial_J(u)}J\bigr)\xi.
\end{equation}
特别地，当 $J$ 与 $\omega$ \BookRecordIndex{T0348}相容时，
\begin{equation}\label{eq:C-7-11}
  D_u\xi=(\widetilde\nabla\xi)^{0,1}
    +\frac14N_J\bigl(\xi,\partial_J(u)\bigr).
\end{equation}
\end{lemma}
\begin{proof}
注意
\[
  \nabla\xi=\widehat\nabla\xi
    +\frac12J(\nabla_{du}J)\xi
    +\frac14J(\nabla_\xi J)du
    +\frac14(\nabla_{J\xi}J)du,
\]
并且由~\eqref{eq:C-7-3}，
\begin{align*}
  N_J\bigl(\xi,\partial_J(u)\bigr)
  &=J(\nabla_{\partial_J(u)}J)\xi
    -(\nabla_{J\partial_J(u)}J)\xi\\
  &\quad-J(\nabla_\xi J)\partial_J(u)
    +(\nabla_{J\xi}J)\partial_J(u).
\end{align*}
因此
\begin{align*}
  D_u\xi
  &=\frac12\bigl(\nabla\xi+J\nabla\xi\circ j
            -J(\nabla_\xi J)\partial_J(u)\bigr)\\
  &=\frac12\bigl(\widehat\nabla\xi
            +J(u)\widehat\nabla\xi\circ j\bigr)
            -\frac12J(\nabla_\xi J)\partial_J(u)\\
  &\quad+\frac14J(\nabla_{du}J)\xi
    +\frac18J(\nabla_\xi J)du
    +\frac18(\nabla_{J\xi}J)du\\
  &\quad-\frac14(\nabla_{du\circ j}J)\xi
    -\frac18(\nabla_\xi J)du\circ j
    +\frac18J(\nabla_{J\xi}J)du\circ j\\
  &=(\widehat\nabla\xi)^{0,1}
    +\frac14J(\nabla_{du}J)\xi
    -\frac14(\nabla_{du\circ j}J)\xi\\
  &\quad-\frac14J(\nabla_\xi J)\partial_J(u)
    +\frac14(\nabla_{J\xi}J)\partial_J(u)\\
  &=(\widehat\nabla\xi)^{0,1}
    +\frac14N_J\bigl(\xi,\partial_J(u)\bigr)
    +\frac14J(\nabla_{\bar\partial_J(u)}J)\xi
    -\frac14(\nabla_{\bar\partial_J(u)\circ j}J)\xi.
\end{align*}
这证明了~\eqref{eq:C-7-10}。公式~\eqref{eq:C-7-11} 由~\eqref{eq:C-7-10} 与~\eqref{eq:C-7-5} 得到。
\end{proof}
